\section{Conclusion}

We asked whether static analyzers can improve LLM code repair. The answer remains open---but we now have the methodology to find it.

\textbf{What we found:} Canonical HumanEval solutions exhibit only 9.15\% pylint warning rates. This makes them invalid proxies for LLM-generated code. Our MUST\_WORK gate correctly identified this limitation before downstream experiments wasted resources.

\textbf{What we built:} A validated pylint analysis pipeline for code generation research, reusable components for HumanEval loading and static analysis, and a gate-based hypothesis validation framework.

\textbf{What we learned:} Gate-based validation works. By testing existence assumptions before full experiments, we catch methodology issues early.

\textbf{The path forward:} Re-run the existence gate with actual LLM generations. If warning rates exceed 30\%, proceed to mechanism validation and full pass@k comparison.

\textbf{Broader implication:} Before evaluating any LLM intervention, validate that the intervention has signal. Canonical solutions are not LLM output---a lesson applicable beyond static analysis to any benchmark-based evaluation.

We contribute methodology, not conclusions. The question of whether static analysis improves pass@k is worth answering---and we provide the tools to answer it.
